% GENERATED by notebooks/scripts/generate_param_register.py from
% simulation/config.py -- do not edit by hand; re-run the script instead.
\begingroup
\setlength{\tabcolsep}{3pt}
\newcolumntype{R}[1]{>{\scriptsize\raggedright\arraybackslash}p{#1}}
\begin{longtable}{R{2.5cm}R{1.1cm}R{1.75cm}R{1.1cm}R{3.8cm}R{2.9cm}}
\caption{The full parameter register}
\label{tab:param-register} \\
\multicolumn{6}{p{13.4cm}}{\footnotesize Every model parameter, generated from
the parameter declarations in \texttt{simulation/config.py} so that the table and
the code cannot drift apart. Columns: the parameter name in the code; the decision
rule or submodel it implements (D$n$ = Submodel $n$; P$n$ = data-layer phase);
its baseline value; its provenance; the source fixing or motivating the value; and
the sweep range where one applies. Provenance: sourced = a citation fixes the
value; derived = the value is computed from data, not chosen; assumed = no anchor
exists and sensitivity analysis is mandatory. Rand values are 2017 Rands; the
\ac{BNPL} contract parameters are published at 2026 vintage and deflated once
(Appendix~\ref{app:variables}, Part~D.12).} \\
\toprule
\textbf{Parameter} & \textbf{Rule} & \textbf{Baseline} & \textbf{Prov.} & \textbf{Source} & \textbf{Sweep} \\
\midrule
\endfirsthead
\toprule
\textbf{Parameter} & \textbf{Rule} & \textbf{Baseline} & \textbf{Prov.} & \textbf{Source} & \textbf{Sweep} \\
\midrule
\endhead
\bottomrule
\endfoot
\texttt{n\_\allowbreak{}ticks} & ODD & 52 & sourced & 24-month horizon at a 14-day tick (03\_calibration.tex) & -- \\
\addlinespace
\texttt{burn\_\allowbreak{}in} & ODD & 12 & sourced & First 12 ticks discarded (03\_calibration.tex) & -- \\
\addlinespace
\texttt{n\_\allowbreak{}agents} & P3 & None & derived & None = use the full 5,000-agent resample & 1000 / 5000 / 10000 population-stability check \\
\addlinespace
\texttt{shock\_\allowbreak{}prob} & D1 & 0.01 & assumed & CALIBRATION TARGET, fitted to CCMR baseline arrears. Reported against the QLFS 2017 job-separation band (0.54\%-1.17\% per tick), NOT fitted to it. & calibration grid \\
\addlinespace
\texttt{shock\_\allowbreak{}persistent} & D1 & True & sourced & An unemployment spell persists until re-employment. Madeira models FLOWS into and out of unemployment, and QLFS 2017 shows 68.4\% of the unemployed remain unemployed the next quarter. & False = the original non-persistent rule, reported as a robustness arm \\
\addlinespace
\texttt{shock\_\allowbreak{}exit\_\allowbreak{}prob} & D1 & 0.018724 & sourced & Per-tick re-employment hazard. QLFS 2017 Q3$\to$Q4: 11.6\% of the unemployed moved into employment per quarter $\to$ 1.87\% per 14-day tick. & with the QLFS band \\
\addlinespace
\texttt{discretionary\_\allowbreak{}floor} & D2 & 0.0 & sourced & Discretionary spend fully compressible in the baseline (D2) & 0.25 / 0.50 habit-persistence floors \\
\addlinespace
\texttt{q\_\allowbreak{}base} & D3/\allowbreak{}D17 & 0.05 & assumed & Spontaneous want-driven BNPL propensity. NOT SOURCED; swept. & RQ2 adoption sweep; sets the level that beta then amplifies \\
\addlinespace
\texttt{beta} & D17 & 0.0 & assumed & Peer imitation strength under the LINEAR mechanism. NOT SOURCED. & RQ2 primary experimental axis -- it is what makes peer effects in BNPL adoption visible -- and toggled on/off for RQ3, whose cool-off lever only bites where beta \textgreater{} 0. The beta=0 row is always reported. \\
\addlinespace
\texttt{peer\_\allowbreak{}mechanism} & D17 & linear & sourced & Which social-transmission rule is active. 'linear' is q = q\_base + beta*s\_g. & RQ2: both mechanisms on identical access grids \\
\addlinespace
\texttt{mu\_\allowbreak{}theta} & D17 & 0.3 & assumed & Mean adoption threshold: the share of a household's reference group that must already be using BNPL before it responds. & 3-point sensitivity at one access level, not a full cross \\
\addlinespace
\texttt{sigma\_\allowbreak{}theta} & D17 & 0.2 & sourced & Dispersion of adoption thresholds, truncated to [0,1]. THE PRIMARY EXPERIMENTAL AXIS of the threshold arm, because Granovetter's actual claim is that the VARIANCE of thresholds, not their mean, decides whether a cascade occurs: a chain of thresholds with someone standing at every level propagates, a tightly clustered one does not. & MANDATORY: 0.05-0.40 \\
\addlinespace
\texttt{gamma} & D17 & 0.0 & assumed & Appetite increment once a household's threshold is crossed. NOT SOURCED. & RQ2 threshold arm; the gamma=0 row is always reported \\
\addlinespace
\texttt{amount\_\allowbreak{}rule} & D4 & shortfall & assumed & No anchor found in the literature sweep. The model's first uncited rule. & MANDATORY: shortfall / +25\% / plus one tick of committed expenditure \\
\addlinespace
\texttt{shortfall\_\allowbreak{}bnpl\_\allowbreak{}capped} & D5 & True & assumed & Caps the BNPL share of a SHORTFALL request at kappa x the household's own monthly budget -- the same quantity that sizes a want-driven purchase. & Robustness: uncapped, crossed with the three D4 amount rules \\
\addlinespace
\texttt{bnpl\_\allowbreak{}purchase\_\allowbreak{}base} & D4/\allowbreak{}D11 & discretionary & sourced & What a BNPL purchase is sized AGAINST. BNPL finances discretionary consumption, so the household's own discretionary budget is the observed scale at which it makes discretionary purchases. & MANDATORY: discretionary / income \\
\addlinespace
\texttt{bnpl\_\allowbreak{}purchase\_\allowbreak{}ratio} & D4/\allowbreak{}D11 & 0.1387 & derived & kappa. Share of a household's discretionary budget spent on the categories BNPL finances (clothing and footwear, furniture and appliances, ICT devices, recreational durables), DERIVED from Stats SA IES 2022/23 COICOP microdata -- see ies\_2022\_bnpl\_share.json. & MANDATORY: 0.07-0.28 (half to double the derived share) \\
\addlinespace
\texttt{bnpl\_\allowbreak{}purchase\_\allowbreak{}cv} & D4 & 0.6 & assumed & Dispersion of BNPL purchase size around its household-relative mean. & with bnpl\_purchase\_ratio \\
\addlinespace
\texttt{min\_\allowbreak{}payer\_\allowbreak{}share} & D6 & 0.29 & sourced & Keys \& Wang (2019): 29\% of accounts pay at or near the minimum. & 0.20-0.40 (brackets the transferred US point estimate) \\
\addlinespace
\texttt{payment\_\allowbreak{}friction} & D6 & 0.0 & sourced & Per-tick probability of missing a traditional instalment DESPITE having the cash. & calibration grid; second fitted parameter \\
\addlinespace
\texttt{min\_\allowbreak{}payment\_\allowbreak{}frac} & D6 & 0.05 & assumed & The contractual minimum itself. D6 fixes the SHARE of minimum-payers but never defined the minimum; the NCA prescribes no formula either. & MANDATORY: 0.025-0.10 \\
\addlinespace
\texttt{k\_\allowbreak{}default} & D7 & 7 & sourced & 7 ticks = 98 days, inside the CCMR 91-120 band, so it maps onto the 90+ impairment convention. & k=4 (56 days $\to$ the 60+ band, target 16.54\%) \\
\addlinespace
\texttt{bnpl\_\allowbreak{}bureau\_\allowbreak{}visible} & D10/\allowbreak{}D14 & False & sourced & BNPL sits outside the NCA: no bureau reporting obligation (Norton Rose; TransUnion). & RQ3 lever 1 - the comparison the model was built to make \\
\addlinespace
\texttt{bnpl\_\allowbreak{}enabled} & O.1a & False & sourced & Injection design: the 2017 baseline is BNPL-free by construction. & the injection itself \\
\addlinespace
\texttt{bnpl\_\allowbreak{}access\_\allowbreak{}rate} & D11 & 1.0 & derived & Share of the BANKED subpopulation that is BNPL-eligible. Banked = 82.8\% of households by observed data, so access\_rate=1.0 means 82.8\% of all households. & RQ2 axis 1, 0.0-1.0 within the banked subpopulation \\
\addlinespace
\texttt{n\_\allowbreak{}platforms} & D12 & 4 & sourced & PayJustNow, Payflex, Mobicred, TymeBank all active in South Africa. & 1-6; N=1 isolates single-platform accumulation from cross-firm stacking \\
\addlinespace
\texttt{bnpl\_\allowbreak{}order\_\allowbreak{}cap} & D11 & 9926.83 & sourced & Payflex per-order cap of R15,000 as published (2026 vintage), DEFLATED to 2017 Rands (cpi\_deflator\_2017.json). & verify it binds rarely at LMI incomes \\
\addlinespace
\texttt{bnpl\_\allowbreak{}limit\_\allowbreak{}income\_\allowbreak{}multiple} & D11 & 0.1 & assumed & lambda. Rolling available balance per platform, as a multiple of the household's MONTHLY income, applied independently at each of n\_platforms. & MANDATORY: 0.1-1.0, with the binding rate reported BY QUINTILE \\
\addlinespace
\texttt{bnpl\_\allowbreak{}instalments} & D13 & 4 & sourced & Payflex Pay in 4: 25\% at checkout then 25\% at each of the next three ticks. & Pay in 3 monthly (PayJustNow) as the structural alternative \\
\addlinespace
\texttt{bnpl\_\allowbreak{}late\_\allowbreak{}fee\_\allowbreak{}per\_\allowbreak{}tick} & D13 & 125.74 & sourced & Payflex late fee R95 per week = R190 per 14-day tick as published (2026 vintage), DEFLATED to 2017 Rands (DEFECTS.md B29). & -- \\
\addlinespace
\texttt{bnpl\_\allowbreak{}late\_\allowbreak{}fee\_\allowbreak{}cap} & D13 & 188.61 & sourced & Payflex caps the late fee at three weeks: 3 x R95 = R285 per missed instalment as published (2026 vintage), DEFLATED to 2017 Rands (DEFECTS.md B29). & -- \\
\addlinespace
\texttt{bnpl\_\allowbreak{}affordability\_\allowbreak{}check} & D14 & False & sourced & Lever 2. Applies the D9 Reg 23A residual-income test to BNPL. & RQ3 lever 2, on/off \\
\addlinespace
\texttt{k\_\allowbreak{}cool} & D14 & 0 & sourced & Lever 3. Baseline 0 = lever OFF. k\_cool=1 tick = 14 days = the UK CCA s.66A statutory right of withdrawal, which lands exactly on a tick boundary. & RQ3 lever 3, 0-4 ticks \\
\addlinespace
\texttt{stacking\_\allowbreak{}cap} & D14 & None & assumed & Lever 4. Max concurrent facilities. HYPOTHETICAL: no jurisdiction imposes one. & RQ3 lever 4, 1-4 or None \\
\addlinespace
\texttt{activation} & D16 & random & sourced & Random asynchronous order, re-drawn every tick (Comer \& Loerch; Alizadeh). & MANDATORY: uniform and synchronous robustness runs \\
\end{longtable}
\endgroup
